\chapter{Mapa dos conceitos que o projeto vai treinar}

Este capitulo apresenta cada conceito rapidamente e mostra onde ele aparece no projeto. A ideia e voce reconhecer o nome, entender a intuicao e depois aplicar no codigo.

\section{I/O-bound}

I/O significa \textit{Input/Output}. Um trecho e I/O-bound quando o programa passa boa parte do tempo esperando algo externo a CPU.

Exemplos:

\begin{itemize}
  \item ler arquivo do disco;
  \item escrever arquivo;
  \item inserir dados no Postgres;
  \item consultar banco;
  \item fazer request HTTP;
  \item esperar uma resposta de rede.
\end{itemize}

No projeto:

\begin{verbatim}
ler export.xml          -> I/O
inserir no Postgres     -> I/O
ler arquivos .gpx       -> I/O
salvar logs             -> I/O
\end{verbatim}

\section{CPU-bound}

Um trecho e CPU-bound quando o gargalo e computacao. O programa esta gastando tempo calculando, convertendo ou processando dados.

No projeto:

\begin{verbatim}
parsear XML em objetos Python       -> parcialmente CPU-bound
converter strings para datas        -> CPU-bound leve
calcular features e estatisticas    -> CPU-bound
normalizar milhoes de pontos GPS    -> CPU-bound
\end{verbatim}

\section{GIL}

GIL significa \textit{Global Interpreter Lock}. No CPython, ele impede que varias threads executem bytecode Python puro ao mesmo tempo no mesmo processo.

Consequencia pratica:

\begin{itemize}
  \item Threads ajudam bastante quando o gargalo e I/O.
  \item Threads geralmente nao aceleram muito trabalho CPU-bound puro.
  \item Multiprocessing ajuda CPU-bound porque cada processo tem seu proprio interpretador e seu proprio GIL.
  \item Async ajuda I/O concorrente, mas nao deixa calculo CPU-bound magicamente paralelo.
\end{itemize}

\section{Stream}

Stream e processar dados aos poucos. Em vez de carregar tudo, voce consome item por item.

\begin{lstlisting}[style=devbutterpython]
for record in stream_records(xml_file):
    process(record)
\end{lstlisting}

No projeto, \code{ET.iterparse} sera a base do streaming.

\section{Fila}

Fila e uma estrutura FIFO: \textit{first in, first out}. O primeiro item que entra e o primeiro que sai.

No projeto:

\begin{verbatim}
producer/parser -> Queue -> consumer/writer Postgres
\end{verbatim}

A fila desacopla velocidades diferentes. Se o parser for mais rapido que o banco, a fila absorve temporariamente. Se a fila encher, o parser espera. Isso e \textbf{backpressure}.

\section{Batching}

Batching e agrupar varias operacoes em uma so.

Ruim:

\begin{verbatim}
1 registro -> INSERT
1 registro -> INSERT
1 registro -> INSERT
\end{verbatim}

Melhor:

\begin{verbatim}
5000 registros -> INSERT batch
5000 registros -> INSERT batch
5000 registros -> INSERT batch
\end{verbatim}

O algoritmo ainda e O(n), mas o fator constante melhora muito.

\section{Logging}

Logging e registrar eventos importantes do programa. Neste projeto, logs devem mostrar:

\begin{itemize}
  \item quantos registros foram processados;
  \item tamanho do batch;
  \item tamanho da fila;
  \item tempo por insert;
  \item registros por segundo;
  \item erros de parsing ou banco.
\end{itemize}

\section{Profiling}

Profiling e medir onde o programa gasta tempo. Sem profiling, otimizacao vira chute.

Medidas importantes:

\begin{itemize}
  \item tempo total;
  \item tempo por batch;
  \item registros por segundo;
  \item uso de memoria;
  \item query time no banco.
\end{itemize}

\section{Big O aplicado}

\begin{longtable}{p{3cm}p{10cm}}
\toprule
\textbf{Complexidade} & \textbf{Aplicacao no projeto} \\
\midrule
O(1) & acesso medio em \code{dict}, \code{set}, append em lista \\
O(n) & percorrer todos os elementos do XML uma vez \\
O(log n) & busca usando indice B-tree no Postgres \\
O(n log n) & ordenacao de registros por data \\
O(n²) & comparacoes ingênuas entre todos os pares de duas listas; algo a evitar \\
\bottomrule
\end{longtable}

\section{Resumo do mapa}

\begin{longtable}{p{4cm}p{9cm}}
\toprule
\textbf{Conceito} & \textbf{Onde aparece} \\
\midrule
list / array & batches de registros \\
dict / hash map & metadata, contadores, mapeamento de tipos \\
set / hash set & filtro rapido de tipos aceitos \\
queue & producer/consumer \\
tree & XML e uma arvore \\
stack & percursos DFS opcionais \\
recursion & exploracao de XML pequeno; evitar para XML gigante \\
sorting & ordenar por data \\
binary search & achar registros em intervalos ordenados \\
two pointers & cruzar intervalos de sono com batimentos \\
sliding window & medias moveis de 7/14 dias \\
caching & summary tables e materialized views \\
indexing & indices por tipo e data no Postgres \\
\bottomrule
\end{longtable}
